\subsection{2. Hướng phát triển}

\begin{itemize}[label={--}, itemsep=1pt]
    \item \textbf{Cơ chế học tăng cường:} Trong môi trường Web3, các tổ chức tạo Sybil liên tục thay đổi. Việc ứng dụng Học tăng cường sẽ giúp hệ thống có khả năng tự động học giá trị phần thưởng từ các phản hồi của chuyên gia, từ đó tự tinh chỉnh lại mô hình để đối phó với các mô thức tấn công mới mà không cần huấn luyện lại từ đầu.
    
    \item \textbf{Mở rộng mô hình hóa đồ thị không gian - thời gian:} Bổ sung chiều thời gian vào cấu trúc đồ thị hiện tại để theo dõi sự tiến hóa của mạng lưới. Mạng nơ-ron đồ thị thời gian sẽ giúp hệ thống nắm bắt được trình tự động học của chuỗi hành vi, từ đó phát hiện sớm các chiến dịch tấn công Sybil ngay từ giai đoạn chúng mới bắt đầu hình thành liên kết thay vì chờ cấu trúc kịch bản hoàn thiện.
    
    \item \textbf{Khung giải thích đa tầng:} Nâng cấp tính minh bạch của toàn bộ pipeline bằng cách kết hợp sức mạnh của hai trường phái giải thích. Cụ thể, sử dụng \textit{SHAP (SHapley Additive exPlanations)} bao bọc bộ phân lớp Random Forest ở giai đoạn cuối để chỉ ra mức độ đóng góp của từng đặc trưng ẩn. Sự kết hợp này sẽ cung cấp bằng chứng suy luận toàn diện và thuyết phục.
\end{itemize}